\section*{Hoofdstuk \ref{ch:b2:himalaya-andes} --- Botsingsgordels en Andesgordels vergeleken}

\begin{solution}{exo:b2:himalaya-andes:1}
(a) $r = h\rho_c/(\rho_m - \rho_c) = 3.0\times2700/600 = \qty{13.5}{km}$. (b) $T = 35 + 3.0 +
13.5 = \qty{51.5}{km}$, of $35 + 5.5\times3.0$. (c) $h = \Delta T(\rho_m - \rho_c)/\rho_m =
10\times600/3300 = \qty{1.8}{km}$.
\end{solution}

\begin{solution}{exo:b2:himalaya-andes:2}
(a) Indo-Gangetische voordiepte; Frontale Hoofdoverschuiving; sub-Himalaya (Siwaliks);
Hoofdgrensoverschuiving; Lage Himalaya; Centrale Hoofdoverschuiving; Hoge-Himalayareeks; Zuid-Tibetaans
Onthechtingssysteem; Tethys-Himalayareeks; daarna de Indus--Tsangpo-sutuur en de Aziatische
korst. (b) Het Zuid-Tibetaanse Onthechtingssysteem. (c) De graad stijgt naar het noorden: niet-gemetamorfoseerde
\omterm{def:b2:alpine-orogen:foreland}{molasse}, middelhooggradige Lage Himalaya, hooggradige en deels gesmolten
Hoge-Himalayareeks; boven het STDS daalt hij abrupt tot niet of zwak gemetamorfoseerde
Tethyssedimenten, omdat de afschuiving de gesteenten ertussen heeft weggenomen.
\end{solution}

\begin{solution}{exo:b2:himalaya-andes:3}
(a) $196.5\ \text{m}/9200\ \text{jaar} = \qty{0.02136}{m}$ per jaar, ongeveer \qty{21.4}{mm} per jaar.
(b) Het is \qty{20}{\percent} meer dan de \qty{17.8}{mm} per jaar van GPS, maar van dezelfde grootte:
de ene snelheid is uitgemiddeld over negenduizend jaar, de andere over twintig.
(c) Bijna alle verkorting over de keten wordt aan het front opgevangen: de MBT, de MCT en
de andere breuken ten noorden van de MFT zijn op deze lengtegraad praktisch inactief; ze wortelen alle
in dezelfde MHT, die bij de MFT aan het oppervlak komt.
\end{solution}

\begin{solution}{exo:b2:himalaya-andes:4}
(a) Botsing: de sutuur markeert de verdwenen oceaan. (b) Andestype: een boog vraagt een plaat
die naar beneden gaat. (c) Botsing, typisch: korstsmelten van een verdikte, verhitte korst (ze komen ook
voor in verdikte Andeskorst, maar daar overheersen de boogmagma's). (d) Andestype.
(e) Beide: de Siwaliks en het Subandiene voorland. (f) Beide: de MHT in de Himalaya, het
subductievlak in de Andes.
\end{solution}

\begin{solution}{exo:b2:himalaya-andes:5}
(a) $17.8/45 = \qty{40}{\percent}$ in het oosten, $20.5/45 = \qty{46}{\percent}$ in het westen:
ongeveer de helft. (b) Naar de verdikking en de breuken van Tibet en de rest van Azië, en naar
de extrusie van blokken naar het oosten langs de grote horizontale verschuivingen. (c) $45 - 17.8 =
\qty{27.2}{mm}$ per jaar, of \qty{272}{km} in \qty{10}{Myr}; $45 - 20.5 = \qty{24.5}{mm}$ per
jaar, of \qty{245}{km}.
\end{solution}

\begin{solution}{exo:b2:himalaya-andes:6}
(a) $r = 4.5\times3.7 = \qty{16.65}{km}$; $\Delta F = 2700\times9.8\times3700\times(35\,000 +
1850 + 8325) = \qty{4.42e12}{N.m^{-1}}$. (b) Tibet: \qty{6.12e12}{N.m^{-1}}; de Altiplano
heeft er \qty{72}{\percent} van. (c) Met $r = 4.5h$ is $\Delta F = \rho_c g\,(Hh + 2.75h^2)$: voor
$h \ll H$ overheerst de term $Hh$ en is $\Delta F$ evenredig met $h$; voor een hoog plateau
groeit de term in $h^2$ sneller. Elke extra kilometer hoogte kost steeds meer duw aan
de randen, tot de beschikbare krachten het plateau niet meer kunnen opheffen; het vloeit dan uit
in plaats van te rijzen.
\end{solution}

\begin{solution}{exo:b2:himalaya-andes:7}
(a) $2700\times9.8\times8848.86/2 = \qty{1.17e8}{Pa} = \qty{117}{MPa}$. (b) Bij
\qty{100}{MPa}: $0.85\times100 = \qty{85}{MPa}$; bij \qty{200}{MPa}: $0.85\times200 = 50 +
0.6\times200 = \qty{170}{MPa}$, de twee takken komen overeen. (c) $h = 2\sigma_d/(\rho g) =
2\times1.70\times10^8/(2700\times9.8) = \qty{12.8}{km}$: een orde van grootte, van de orde
van de hoogste toppen, er niet ver boven.
\end{solution}

\begin{solution}{exo:b2:himalaya-andes:8}
(a) $11\times111.19\times\cos\ang{30} = \qty{1059}{km}$. (b) $9\times10^{-3}/1.059\times10^{6}
= \num{8.5e-9}$ per jaar (tussen \num{6.6e-9} en \num{10.4e-9} voor 7 en \qty{11}{mm} per jaar),
dat wil zeggen $8.5\times10^{-9}/3.156\times10^{7} = \qty{2.7e-16}{s^{-1}}$. (c) De twee bereiken
overlappen: $(12.4 \pm 4.7) \times 10^{-9}$ loopt van 7.7 tot \num{17.1e-9}. Het zuidelijke plateau ten noorden van
Nepal rekt even snel uit als de gordel als geheel, of iets sneller.
\end{solution}

\begin{solution}{exo:b2:himalaya-andes:9}
(a) $0.9/6.6 = \qty{14}{\percent}$. (b) $M_0 = 10^{1.5\times8.5 + 9.1} = 10^{21.85} =
\qty{7.1e21}{N.m}$. (c) $7.1\times10^{21}/6.6\times10^{19} = 107$: één zo'n aardbeving ongeveer
om de 110 jaar onder Nepal, veel meer dan de geschiedenis optekent; ofwel staan er veel grotere aardbevingen
te gebeuren, ofwel komt een deel van het tekort op een andere manier vrij.
\end{solution}

\begin{solution}{exo:b2:himalaya-andes:10}
(a) Neem de compensatiediepte aan de basis van de korst van de Altiplano, $61.5 - 3.7 =
\qty{57.8}{km}$ onder zeeniveau. De korst van de Puna reikt tot $50.5 - 4.4 = \qty{46.1}{km}$, dus de
kolom van de Puna bevat boven die diepte \qty{11.7}{km} mantel: $2800\times61.5 =
\rho_P\times50.5 + 3300\times11.7$. (b) $\rho_P = (172\,200 - 38\,610)/50.5 =
\qty{2645}{kg.m^{-3}}$. (c) Het verschil, \qty{155}{kg.m^{-3}}, ligt aan de onderkant van
de opgegeven 150 tot \qty{200}{kg.m^{-3}}: de korst alleen zou een ongewoon lichte Puna-korst vragen.
Een lichtere kolom eronder doet het gemakkelijker: een mantellithosfeer die door delaminatie is verwijderd,
vervangen door hete, lichte asthenosfeer.
\end{solution}

\begin{solution}{exo:b2:himalaya-andes:11}
(a) $g = 105/\tan\ang{30} = \qty{181.9}{km}$. (b) $g = 90/\tan\ang{30} + 15/\tan\ang{4} =
155.9 + 214.5 = \qty{370.4}{km}$: de boog verschuift \qty{188.5}{km} landinwaarts. (c) De plaat bereikt
over die \qty{300}{km} nooit \qty{105}{km} onder de plaat, en er ligt geen hete mantelwig
boven: er wordt geen \omterm{def:b2:magma-emplacement:magma}{magma} gemaakt en de boog sterft af tot de plaat weer wegzinkt, verder naar het oosten.
\end{solution}

\begin{solution}{exo:b2:himalaya-andes:12}
(a) \qty{21}{mm} per jaar is \qty{21}{km} per miljoen jaar: $21\times37.5 =
\qty{787.5}{km^3}$ per kilometer keten per miljoen jaar. (b) $300/787.5 =
\qty{38}{\percent}$. (c) In de verdikking van de keten en van Zuid-Tibet, in korst die
met India onder Tibet wordt meegevoerd, en in het materiaal dat de MHT naar het noorden doorgeeft.
\end{solution}

\begin{solution}{pb:b2:himalaya-andes:1}
\textbf{1.} Zeven monsters van \ang{28.0} tot \ang{37.0}~NB: $34\,550/7 = \qty{4936}{m}$, $h_T
= \qty{4.9}{km}$.
\textbf{2.} Vijf monsters van \ang{-69} tot \ang{-65}: $18\,650/5 = \qty{3730}{m}$, $h_A =
\qty{3.7}{km}$, de waarde die voor de Altiplano wordt opgegeven; de Puna, verder naar het zuiden, ligt hoger.
\textbf{3.} Tibet: minstens $9\times111.19 = \qty{1001}{km}$; de Altiplano: minstens $4\times
104.5 = \qty{418}{km}$.
\textbf{4.} $(5000 - 70)/(1.5\times111\,190) = 0.0296$: \qty{3.0}{\percent}.
\textbf{5.} Het Andesprofiel begint op de oceaanbodem en in de trog, kilometers onder
zeeniveau, ten westen van het plateau; het Tibetaanse begint op een continentale vlakte dicht bij zeeniveau,
en eindigt op het hooggelegen Tarimbekken.
\textbf{6.} $r = 4.5\times4.9 = \qty{22.05}{km}$; $T_T = 35 + 4.9 + 22.05 = \qty{62.0}{km}$.
\textbf{7.} $r = 4.5\times3.7 = \qty{16.65}{km}$; $T_A = 35 + 3.7 + 16.65 = \qty{55.4}{km}$.
\textbf{8.} Tibet: \qty{62.0}{km} tegen 70 tot \qty{75}{km}, \qty{10.5}{km} onder de middelste
waarde \qty{72.5}{km} (\qty{14}{\percent}); de Altiplano: \qty{55.4}{km} tegen 59 tot
\qty{64}{km}, \qty{6.1}{km} onder \qty{61.5}{km} (\qty{10}{\percent}). Tibet zit er verder naast.
\textbf{9.} $72.5 - 35 = 4.9\times3300/(3300 - \rho_c)$: $3300 - \rho_c = 16\,170/37.5 =
431.2$, $\rho_c = \qty{2869}{kg.m^{-3}}$.
\textbf{10.} $26.5 = 3.7\times3300/(3300 - \rho_c)$: $3300 - \rho_c = 460.8$, $\rho_c =
\qty{2839}{kg.m^{-3}}$.
\textbf{11.} Een deel van de last wordt gedragen door het buigen van de plaat, niet door een lokale wortel; de
normale korst was misschien niet \qty{35}{km} dik; de mantel eronder kan lichter of dichter zijn
dan aangenomen; de seismische Moho is misschien niet de basis van de korst in de zin van de dichtheid.
\textbf{12.} De basis van de korst van Tibet ligt op $72.5 - 4.9 = \qty{67.6}{km}$, die van India op
\qty{37.5}{km}: $t = \qty{30.1}{km}$. $\Delta g = 2\pi\times6.674\times10^{-11}\times370\times
30\,100 = \qty{4.67e-3}{m.s^{-2}}$, dat wil zeggen \qty{-467}{mGal} ten opzichte van India, dicht bij het
waargenomen \qty{460}{mGal}.
\textbf{13.} $W_0 = 440\times60/35 = \qty{754.3}{km}$; $S = \qty{314.3}{km}$.
\textbf{14.} $(314.3 - 250)/314.3 = \qty{20}{\percent}$, het tekort dat in de Centrale
Andes is gevonden.
\textbf{15.} $0.55\times314.3 = \qty{173}{km}$, minder dan de \qty{250}{km} van de \omterm{def:b2:alpine-orogen:balanced}{gebalanceerde
doorsneden}, die dan ruim zouden volstaan.
\textbf{16.} Het dek dat boven het \omterm{def:b2:alpine-orogen:decollement}{décollement} is verkort, moet worden gecompenseerd door verkorting van de
sokkel eronder, verder naar het westen, omdat beide over de gordel dezelfde lengte moeten verliezen; de verkorting van de sokkel
verdikt de korst onder de Oostelijke Cordillera en het plateau.
\textbf{17.} $314.3/30 = \qty{10.5}{mm}$ per jaar, $10.5/72 = \qty{15}{\percent}$ van de
convergentie.
\textbf{18.} Ze wordt door subductie verbruikt: de Nazcaplaat zinkt onder het continent.
\textbf{19.} $r = \qty{22.05}{km}$; $\Delta F = 2700\times9.8\times4900\times(35\,000 + 2450
+ 11\,025) = \qty{6.28e12}{N.m^{-1}}$.
\textbf{20.} $2700\times9.8\times3700\times45\,175 = \qty{4.42e12}{N.m^{-1}}$; de verhouding is
0.70.
\textbf{21.} $6.28\times10^{12}/10^{5} = \qty{6.3e7}{Pa}$, ongeveer \qty{63}{MPa}.
\textbf{22.} $9\times10^{-3}/1.06\times10^{6} = \num{8.5e-9}$ per jaar.
\textbf{23.} De duw van India werkt noord--zuid en wordt in het hoogplateau in evenwicht gehouden door het
eigen gewicht van het plateau: noord--zuid kan het plateau niet meer verdikken, zodat de verticale spanning,
het gewicht, de grootste is; de kleinste spanning is oost--west, naar de vrije randen van Azië,
en het plateau rekt in die richting uit, langs noord--zuid-riften.
\textbf{24.} \textbf{\qty[text-series-to-math=true]{62}{km}}, ongeveer \qty{86}{\percent} van de
gemeten \qty{72.5}{km}.
\end{solution}
